\section{Experiment}
\label{s:experiment}
\subsection{Data Description}


In this study, we analyzed 111 patients who were diagnosed with advanced ovarian cancer between March 2018 and June 2020, and whose diagnosis was confirmed through pathology. All patients underwent cytoreduction surgery performed by the same surgical team. As shown in Table~\ref{tab:dataset}, the surgical results for 71 patients were R0, while for the remaining 40 patients, they were non-R0. We randomly selected 88 patients out of the 111 (80\%) for the training set, and 23 out of the 111 (20\%) for the test set. Before cytoreduction surgery, each patient received an enhanced CT scan of the abdominal cavity. Two experienced doctors selected a series of CT slices that contained the entire region of the disease for each patient. The number of CT slices in each sample ranged from 16 to 136, with an average of 63.53 and a median of 61, and the height and width of all CT slices were 512.
The above dataset is desensitized and collected with the consent of all patients, and the study has been approved by the ethics review committee.

\begin{table}
 \centering
 \caption{Number of samples (slices) in the dataset.}
 \label{tab:dataset}
 \tabcolsep=3pt
 \scalebox{1.1}{
 \begin{tabular}{lccc}
  \toprule
  \  & Training set & Test set & In total\\
  \midrule
    non-R0 & 32 (1,712) & 8 (371) & 40 (2,083) \\
    R0 & 56 (3,776) & 15 (1,193) & 71 (4,969) \\
    In total & 88 (5,488) & 23 (1,564) & 111 (7,052) \\
  \bottomrule
 \end{tabular}
 }
\end{table}

\begin{figure*}[t]
\centering
\begin{center}
\begin{minipage}{16cm}
\begin{tabular*}{16cm}{@{\extracolsep{\fill}}c@{}c@{}c@{\extracolsep{\fill}}}
\includegraphics[height=45mm, width=50mm]{figure/roc_ml.pdf} &
\includegraphics[height=45mm, width=50mm]{figure/roc_dl.pdf} &
\includegraphics[height=45mm, width=50mm]{figure/roc_ablation.pdf} \\
(a) Machine learning methods & (b) deep-learning methods & (c) Ablation performance\\
\end{tabular*}
\end{minipage}
\end{center}
% \vspace{-4mm}
\caption{ROC curves of different methods.}\label{f:roc}
% \vspace{-4mm}
\end{figure*}

\subsection{Experiment Setting}
\begin{itemize}

\item \textbf{Model Setting:} 
In the CT images, the value of each pixel is an integer called CT number, which reflects the attenuation of the X-ray. 
To highlight the ovaries we truncate the CT number to the $[-100, 200]$ interval. 
After truncation, a max-min normalization is applied to map the CT number to $[0, 1]$ interval. 
Considering the mean and median of the number of CT slices, and the computational complexity, we employ bi-linear interpolation to resize all the depth to 64, and the height and width to 256 of all 3D CT images, i.e. $d = 64, h = w = 256$, the dimensionality-reduction ratio r is set to $8$. 
In the classification part, the output dimension of the backbone network $n$ is set to $512$.

\item \textbf{Training Setting:}  
To deal with the small-scale dataset, we apply data augmentation techniques to process the training set, including random translation and rotation, random flipping, and adding Gaussian noise.
In addition to the backbone using pre-trained weights, we set all weights following the Gaussian distribution to initialize the MuVAL model. 
To fine-tune the model, we apply the adamW \cite{loshchilov2017decoupled} to optimize the loss function, and we further initialize the learning rate $lr=0.0001$ and adjust by the exponential decay with $\gamma=0.99$. Moreover, we set the batch size to $10$, the epoch of training to $60$, and utilize an early stop with the patience of $40$. All experiments are implemented on Python 3.8.10, sklearn 1.0.2, Pytorch 1.11.0, and an NVIDIA RTX A5000 GPU with 24GB memory.

\begin{table}
 \centering
 \caption{Results and the number of parameters of models.}
 \label{tab:result}
  %\resizebox{8.5cm}{!}{
 \tabcolsep=3pt
 \begin{tabular}{lcccc}
  \toprule
Model & \#params & ACC(\%) $\uparrow$ & F1(\%) $\uparrow$ & AUC $\uparrow$\\
  \midrule
    (1)LR & / &  56.52 & 64.29 & 0.5500\\ %
    (2)kNN \cite{cover1967nearest} & / & 52.17 & 68.57 & 0.3708\\ %
    (3)SVM \cite{cortes1995support} &  / &47.83 & 53.85 & 0.5167\\ %
  \midrule
    (4)ResNet-2D \cite{he2016deep} & 21.47 M & 73.91 & 82.35 & 0.6000 \\ %
    (5)ResNet-3D \cite{he2016deep} & 63.30 M & 73.91 & 75.00 & 0.7000\\ %
    (6)SE-ResNet-3D \cite{hu2018squeeze} & 63.62 M & 78.26 & 82.76 & 0.6917\\
  \midrule
    (7)Med3D \cite{chen2019med3d} & 63.30 M & 78.26 & 82.76 & 0.7250\\
    \makecell[l]{(8)MuVAL\\(w/o pre-trained)} & 63.33 M & 78.26 & 85.71 & 0.7917\\
    \makecell[l]{(9)MuVAL\\(single-view)} & 63.31 M & 82.61 & 87.50 & 0.7792\\ %
    (10)MuVAL & 63.33 M & \textbf{86.96} & \textbf{89.66} & \textbf{0.8417}\\ %
  \bottomrule
 \end{tabular}
\end{table}

\item \textbf{Evaluation Metrics:}
    In general, the Accuracy (ACC) and F1-score (F1) are commonly applied evaluation metrics for classification.
    %due to difficult
    However, due to the small-scale test set in our task, the ACC may be difficult to compare models with similar performance. 
    In the medical field, the Receiver Operating Characteristic (ROC) curve and Area Under the receiver operating characteristic Curve (AUC) are usually seen as important metrics, which reflect the inter-class distinction and allow comparison of models with similar ACC. 
    Furthermore, ROC and AUC can evaluate the generalization performance of the model in a certain. 
    In conclusion, we combined the above metrics to evaluate the model.
\end{itemize}

\begin{table}
 \centering
 \caption{Case study of MuVAL.}
 \label{tab:case}
 \tabcolsep=3pt
 \begin{tabular}{lccc}
  \toprule
  Sample & Truth & Model prediction & Probability of R0\\
  \midrule
    (1) & R0 & R0 & 88.72\%\\ %p1 1
    (2) & R0 & R0 & 84.13\%\\ %p2 2
    (3) & R0 & R0 & 86.43\%\\ %p3 3
  \midrule
    (4) & non-R0 & non-R0 & 25.68\%\\ %p4 4
    (5) & non-R0 & non-R0 & 10.57\%\\ %p5 5
    (6) & non-R0 & R0 & 72.71\%\\ %p6 6
  \bottomrule
 \end{tabular}
 % \vspace{-4mm}
\end{table}

\subsection{Experiment Results} 

\begin{figure*}[t]
\centering
\begin{center}
\begin{minipage}{17.5cm}
%\begin{tabular*}{17cm}{{@{\extracolsep{\fill}}m{2cm}<{\centering}c@{}c@{}c@{}c@{}c@{}c@{}c@{\extracolsep{\fill}}}}
\begin{tabular*}{17.5cm}{m{1.5cm}<{\centering} m{2.1cm}<{\centering} m{2.1cm}<{\centering} m{2.1cm}<{\centering} m{2.1cm}<{\centering} m{2.1cm}<{\centering} m{2.1cm}<{\centering}}
 & TP sample (1) & TP sample (2) & TP sample (3) & TN sample (4) & TN sample (5) & FP sample (6)\\
Transverse  &
\includegraphics[width = 2.25cm, height=1.5cm]{figure/p1_0_masked.png} &
\includegraphics[width = 2.25cm, height=1.5cm]{figure/p2_0_masked.png} &
\includegraphics[width = 2.25cm, height=1.5cm]{figure/p3_0_masked.png} &
\includegraphics[width = 2.25cm, height=1.5cm]{figure/p4_0_masked.png} &
\includegraphics[width = 2.25cm, height=1.5cm]{figure/p5_0_masked.png} &
\includegraphics[width = 2.25cm, height=1.5cm]{figure/p6_0_masked.png} \\
Coronal &
\includegraphics[width = 2.25cm, height=0.5cm]{figure/p1_1_masked.png} &
\includegraphics[width = 2.25cm, height=0.5cm]{figure/p2_1_masked.png} &
\includegraphics[width = 2.25cm, height=0.5cm]{figure/p3_1_masked.png} &
\includegraphics[width = 2.25cm, height=0.5cm]{figure/p4_1_masked.png} &
\includegraphics[width = 2.25cm, height=0.5cm]{figure/p5_1_masked.png} &
\includegraphics[width = 2.25cm, height=0.5cm]{figure/p6_1_masked.png} \\
Sagittal &
\includegraphics[width = 2.25cm, height=0.5cm]{figure/p1_2_masked.png} &
\includegraphics[width = 2.25cm, height=0.5cm]{figure/p2_2_masked.png} &
\includegraphics[width = 2.25cm, height=0.5cm]{figure/p3_2_masked.png} &
\includegraphics[width = 2.25cm, height=0.5cm]{figure/p4_2_masked.png} &
\includegraphics[width = 2.25cm, height=0.5cm]{figure/p5_2_masked.png} &
\includegraphics[width = 2.25cm, height=0.5cm]{figure/p6_2_masked.png} \\
\end{tabular*}
\end{minipage}
\end{center}
% \vspace{-4mm}
\caption{The multi-view slices of samples in Table~\ref{tab:case}, with the disease region outlined. We can observe that positive samples have larger disease regions, while negative samples have smaller disease regions and multiple disease regions.}\label{f:case}

% \vspace{-4mm}
\end{figure*}

We compare our method against the following baselines, Logistic Regression (LR), k-NearestNeighbor (kNN) \cite{cover1967nearest}, Support Vector Machine (SVM) \cite{cortes1995support}, ResNet-2D and ResNet-3D \cite{he2016deep}, SE-ResNet-3D \cite{hu2018squeeze}. All machine learning methods apply grayscale histograms as input, and all ResNets have a 34-layer structure. The overall performance is summarized in Table~\ref{tab:result}, and the ROC curves are shown in Fig~\ref{f:roc}. We can observe that MuVAL achieves the best performance, especially the high AUC indicates that our method has the better inter-class discriminatory ability.

As with traditional machine learning methods, experiments (1)-(3) fail in the task, with all metrics disappointing. Neither of them could well represent image information, especially 3D image information.

Compared to traditional machine learning methods, deep-learning methods, experiments (4)-(6), achieve a high improvement. Specifically, ResNet-2D and ResNet-3D achieve the same ACC. However, ResNet-2D has higher F1 while ResNet-3D has higher AUC, indicating that ResNet-2D has better classification on positive samples, while ResNet-3D has better inter-class distinction. Theoretically, ResNet-3D is more suitable for processing 3D CT images than ResNet-2D. But ResNet-3D has a larger number of parameters, making training more difficult, especially on small-scale datasets, which may result in a similar performance to ResNet-3D and ResNet-2D. With the addition of the attention mechanism, SE-ResNet-3D achieves higher ACC, similar F1 as ResNet-2D, and similar AUC as ResNet-3D. 

\subsection{Ablation Performance}
As shown in Table~\ref{tab:result}, we conduct an ablation study, i.e., experiments (7)-(9).
To compare the performance of the attention part and the classification part, in experiment (7), we remove the entire attention part, remaining only the Med3D backbone.
Similarly, in experiment (8), we remove the Med3D backbone and employ a ResNet-3D without pre-trained instead. 
Even though the results show that the ACC for both is the same, the MuVAL method is more important than the pre-trained backbone because MuVAL without pre-trained achieves higher F1 and AUC.

In experiment (9), we take out the multi-view mechanism, leaving only the transverse view. 
The next-highest ACC and F1 are attained by the single-view model, indicating the performance boost from fusing the attention mechanism with the pre-trained backbone.
However, the large performance gap with the full MuVAL demonstrates the importance of the multi-view mechanism.

\subsection{Case study}
To further analyze the performance of MuVAL, Table~\ref{tab:case} shows the classification results of some test samples. 
According to the prediction and probability, our method performs well on the true positive (TP) samples (1)-(3) and the true negative (TN) samples (4) and (5) but not on the false positive (FP) sample (6).
The multi-view slices of these samples are shown in Fig.~\ref{f:case}, with the disease region outlined. We can observe that the true positive (TP) samples (1)-(3) have larger disease regions, while true negative (TN) samples (4) and (5) samples have smaller disease regions and multiple disease regions. 
The false positive (FP) sample (6) has a large disease region,  which may lead to misclassification. 
The above observations show that our method achieves classification by representing the 3D CT images in terms of disease shape and size. 
However, there may be latent or hard-to-represent information in the 3D CT images such as disease invasion, metastasis, and other information outside the abdominal cavity.
Finally, the results show that a large disease region is more likely to be R0.
The medical explanation \cite{horvath2013relationship} is that the early-stage disease is larger and more complete, while the late-stage disease rapidly detaches and metastasizes widely, causing more surgical difficulties.

\subsection{Discussion}
Despite the effectiveness of our method, several issues remain to be further explored in the future. 
At first, we use a small-scale real-world dataset during the experiments.
Therefore, our method needs to be further tested and evaluated on large-scale datasets. 
Secondly, our case study shows that 3D CT images may be difficult to represent disease invasion and metastasis.
The following research is to apply multi-modal data to improve the performance further. 
Finally, our method may have the potential to predict the residual disease of other cancers, which requires more experiments to verify.